% Introduction --- derived from the Avance7 "Problem Statement" and EDA
% findings, reorganized as a paper introduction with an explicit thesis
% and contributions.
\section{Introducción}
\label{sec:intro}

La gestión agrícola eficiente requiere conocer con precisión qué se
cultiva, dónde y en qué cantidad. Los métodos tradicionales basados en
la inspección física en campo son costosos, lentos y difíciles de
escalar a grandes territorios. Las imágenes Sentinel-2 ofrecen una
cobertura continua y de acceso libre, pero convertirlas en mapas de
cultivo confiables es un problema difícil de aprendizaje automático por
tres razones estructurales.

\paragraph{Desbalance de clases severo.}
De las 18 clases de cultivo en PASTIS-R~\cite{garnot2021pastis}, el
prado representa hasta el 30\% de todas las parcelas, mientras que
clases de valor agronómico como la colza, el sorgo y el huerto son
clases minoritarias (por debajo del 1\%). Un modelo que ignora las
clases raras tiene poco valor comercial, de modo que las métricas
promediadas en macro son el objetivo correcto.

\paragraph{Ambigüedad espectral intra-clase.}
Cultivos como el trigo blando de invierno y la cebada de invierno son
prácticamente indistinguibles en una sola imagen RGB. La discriminación
requiere explotar la \emph{dinámica temporal} de la firma espectral a lo
largo del año --- inicio del reverdecimiento, pico de NDVI, fecha de
cosecha ---, lo que motiva las arquitecturas de series temporales frente
a los clasificadores de imagen
única~\cite{tarasiou2023tsvit,garnot2021utae}.

\paragraph{Restricciones de producción.}
Un sistema desplegable debe producir predicciones a nivel de parcela (no
solo a nivel de píxel), ser explicable a usuarios no técnicos y operar
con latencias de inferencia compatibles con una interfaz interactiva.

\paragraph{Tesis.}
Argumentamos que la cuantificación del área de cultivo se atiende mejor
con una arquitectura de \emph{perceptor/razonador}. Una pila de
percepción de modelos especializados --- transformadores temporales,
clasificadores tabulares sobre los embeddings anuales de
AlphaEarth~\cite{brown2025alphaearth} y una rama contrastiva
visión-lenguaje~\cite{li2025farslip} --- produce observaciones
auditables en texto plano por parcela. Un LLM frontera \emph{congelado}
razona luego sobre esas observaciones y se comunica con el usuario,
siguiendo el patrón \emph{Be My
Eyes}~\cite{huang2025bemyeyes}. La capa de lenguaje es un componente de
comunicación, explicación y orquestación de herramientas --- no un
clasificador de píxeles ---, lo que mantiene cada cifra citada trazable
a un resultado de herramienta determinista y preserva la opción de
operación on-premises para la soberanía de datos.

\paragraph{Contribuciones.}
\begin{enumerate}[leftmargin=*]
    \item Un estudio comparativo de arquitecturas de segmentación densa
    sobre PASTIS-R, que identifica a
    TSViT-pheno~\cite{tarasiou2023tsvit} como el mejor modelo individual
    (mIoU 0.6253, F1 macro 0.7500 en el fold de prueba retenido), y
    ensambles heterogéneos que mejoran el F1 macro out-of-fold sobre
    cualquier miembro individual
    (Sección~\ref{sec:results}).
    \item Una capa conversacional anclada construida sobre el patrón
    \emph{Be My Eyes} con un razonador congelado dual (Gemini~3.5~Flash en
    la nube / Qwen3-30B-A3B on-premises) y un conjunto cerrado de
    herramientas geoespaciales con Spatial-RAG
    (Sección~\ref{sec:method}).
    \item Un estudio honesto de la transferibilidad espacial de Francia a
    Cataluña (Sen4AgriNet~\cite{sykas2022sen4agrinet}) y una curva
    few-shot de $k$-shots sobre
    EuroCropsML~\cite{reuss2025eurocropsml}, que muestra que la
    transferencia zero-shot falla y que se requiere la adaptación
    few-shot, en línea con los límites
    reportados~\cite{harvesting2026alphaearth}
    (Sección~\ref{sec:results}).
\end{enumerate}

El conjunto de datos de trabajo es PASTIS-R (Francia metropolitana), con
2{,}433 parches Sentinel-2 multitemporales de $128\times128$ píxeles,
etiquetas reales validadas por agricultores, 18 cultivos y folds
espacialmente disjuntos que evitan la fuga de información entre
entrenamiento y evaluación.
